\documentclass[aspectratio=169,11pt]{beamer}

\usetheme{AnnArbor}
\usecolortheme{dolphin}
\usefonttheme{professionalfonts}
\setbeamertemplate{navigation symbols}{}
\setbeamertemplate{caption}[numbered]
\setbeamertemplate{itemize items}[circle]
\setbeamertemplate{enumerate items}[default]

\usepackage{fontspec}
\setsansfont{DejaVu Sans}
\setmonofont{DejaVu Sans Mono}
\usepackage{amsmath,amssymb}
\usepackage{booktabs}
\usepackage{array}
\usepackage{tabularx}
\usepackage{adjustbox}
\usepackage{listings}
\usepackage{xcolor}
\usepackage{tikz}
\usetikzlibrary{arrows.meta,positioning,shapes.geometric,fit,shadows}
\usepackage{hyperref}
\hypersetup{colorlinks=true,linkcolor=dbblue,urlcolor=dbblue}

% Palette mau thuong hieu bai giang
\definecolor{dbblue}{RGB}{20,74,135}
\definecolor{dblight}{RGB}{235,243,255}
\definecolor{dbgreen}{RGB}{28,120,90}
\definecolor{dborange}{RGB}{210,120,30}
\definecolor{codebg}{RGB}{246,248,250}
\definecolor{codekw}{RGB}{0,80,180}
\definecolor{codestr}{RGB}{180,60,20}
\definecolor{codecomm}{RGB}{100,110,120}

\setbeamercolor{structure}{fg=dbblue}
\setbeamercolor{title}{fg=white,bg=dbblue}
\setbeamercolor{frametitle}{fg=dbblue}
\setbeamercolor{block title}{fg=white,bg=dbblue}
\setbeamercolor{block body}{bg=dblight,fg=black}
\setbeamercolor{block title example}{fg=white,bg=dbgreen}
\setbeamercolor{block body example}{bg=dblight!60!white,fg=black}
\setbeamercolor{block title alerted}{fg=white,bg=dborange}
\setbeamercolor{block body alerted}{bg=dblight!40!white,fg=black}

\lstset{
  basicstyle=\ttfamily\scriptsize,
  keywordstyle=\color{codekw}\bfseries,
  stringstyle=\color{codestr},
  commentstyle=\color{codecomm}\itshape,
  backgroundcolor=\color{codebg},
  frame=single,
  rulecolor=\color{dbblue!30},
  breaklines=true,
  showstringspaces=false,
  keepspaces=true,
  aboveskip=2pt,
  belowskip=2pt,
  language=SQL,
  morekeywords={BOOLEAN,TINYINT,SMALLINT,MEDIUMINT,BIGINT,DECIMAL,NUMERIC,FLOAT,DOUBLE,VARCHAR,CHAR,TEXT,TINYTEXT,MEDIUMTEXT,LONGTEXT,DATE,TIME,DATETIME,TIMESTAMP,YEAR,BINARY,VARBINARY,BLOB,TINYBLOB,MEDIUMBLOB,LONGBLOB,JSON,GEOMETRY,UNSIGNED,ZEROFILL},
  tabsize=2
}

\newcommand{\kw}[1]{\textcolor{dbblue}{\textbf{#1}}}
\newcommand{\smallgap}{\vspace{0.25em}}

\title[Kiểu dữ liệu SQL]{Bài giảng: Kiểu dữ liệu trong SQL \& MySQL}
\subtitle{Phân loại, Lựa chọn tối ưu, Miền giá trị \& Thực hành thiết kế CSDL}
\author[Vũ Đức Minh]{Vũ Đức Minh}
\institute[NEU]{Trường Đại học Kinh tế Quốc dân (NEU), Hà Nội, Việt Nam}
\date{\today}

\begin{document}

%------------------------------------------------
% Slide 1: Trang bia
\begin{frame}
  \titlepage
\end{frame}

%------------------------------------------------
% Slide 2: Muc tieu bai hoc
\begin{frame}{Mục tiêu bài giảng}
Sau khi hoàn thành bài học này, người học có thể:
\begin{itemize}\setlength{\itemsep}{2pt}
  \item Trình bày được khái niệm \kw{kiểu dữ liệu (Data Type)} và vai trò trong RDBMS.
  \item Phân biệt chính xác các nhóm: \kw{Numeric, String, Date/Time, Binary, Boolean \& Special}.
  \item Nắm vững cơ chế lưu trữ của \kw{CHAR vs VARCHAR}, \kw{DECIMAL vs FLOAT}, \kw{DATETIME vs TIMESTAMP}.
  \item Hiểu lý do bắt buộc dùng kiểu số thập phân chính xác (\kw{DECIMAL}) cho tài chính/tiền tệ.
  \item Lựa chọn kiểu dữ liệu tối ưu dung lượng đĩa, bộ nhớ đệm RAM và chỉ mục (indexing).
  \item Vận dụng khai báo kiểu dữ liệu kết hợp ràng buộc trong câu lệnh \kw{CREATE TABLE}.
  \item Nhận biết và khắc phục các lỗi thiết kế kiểu dữ liệu kinh điển trong thực tế.
\end{itemize}
\end{frame}

%------------------------------------------------
% Slide 3: Noi dung bai hoc
\begin{frame}{Nội dung bài học}
\begin{columns}[t]
  \begin{column}{0.48\textwidth}
    \begin{block}{Phần 1: Khái niệm \& Phân nhóm}
      \begin{itemize}\setlength{\itemsep}{2pt}
        \item 1. Bản chất kiểu dữ liệu trong SQL
        \item 2. Sơ đồ các nhóm kiểu dữ liệu chính
        \item 3. Quy trình 5 bước lựa chọn kiểu
        \item 4. Kiểu số nguyên: Range \& Unsigned
        \item 5. Số thực: DECIMAL vs FLOAT/DOUBLE
      \end{itemize}
    \end{block}
  \end{column}
  \begin{column}{0.48\textwidth}
    \begin{block}{Phần 2: Ứng dụng \& Tối ưu}
      \begin{itemize}\setlength{\itemsep}{2pt}
        \item 6. Kiểu chuỗi: CHAR, VARCHAR, TEXT
        \item 7. Kiểu ngày giờ: DATETIME vs TIMESTAMP
        \item 8. Kiểu Binary, Boolean \& JSON
        \item 9. Thực hành thiết kế \& Sai lầm thường gặp
        \item 10. Câu hỏi ôn tập \& Bài tập vận dụng
      \end{itemize}
    \end{block}
  \end{column}
\end{columns}
\end{frame}

%------------------------------------------------
\section{Bản chất kiểu dữ liệu}
% Slide 4: Khai niem kieu du lieu
\begin{frame}{1. Bản chất kiểu dữ liệu trong SQL}
\begin{block}{Định nghĩa}
  \textbf{Kiểu dữ liệu (Data Type)} quy định \kw{loại giá trị}, \kw{miền giá trị (domain)}, \kw{dung lượng bộ nhớ} và \kw{các phép toán hợp lệ} trên một cột.
\end{block}

\smallgap
\begin{columns}[t]
  \begin{column}{0.48\textwidth}
    \textbf{Vai trò cốt lõi trong RDBMS:}
    \begin{itemize}\setlength{\itemsep}{1.5pt}
      \item \kw{Toàn vẹn (Integrity)}: Chặn dữ liệu sai ngay từ tầng lưu trữ.
      \item \kw{Tối ưu hiệu năng}: Giảm I/O đĩa, tăng dung lượng bản ghi trong Buffer Pool.
      \item \kw{Tính toán chuẩn xác}: Đảm bảo phép so sánh, sắp xếp và tính toán không sai lệch.
    \end{itemize}
  \end{column}
  \begin{column}{0.48\textwidth}
    \textbf{Ba thuộc tính then chốt:}
    \begin{itemize}\setlength{\itemsep}{1.5pt}
      \item \kw{Miền giá trị}: Khoảng giá trị hợp lệ (vd: \texttt{TINYINT} $-128 \dots 127$).
      \item \kw{Độ dài (Length)}: Số ký tự tối đa (\texttt{VARCHAR(100)}).
      \item \kw{Độ chính xác (Precision/Scale)}: Số chữ số và phần thập phân (\texttt{DECIMAL(10,2)}).
    \end{itemize}
  \end{column}
\end{columns}
\end{frame}

%------------------------------------------------
% Slide 5: So do phan nhom
\begin{frame}{2. Sơ đồ các nhóm kiểu dữ liệu phổ biến}
\begin{columns}[c]
  \begin{column}{0.48\textwidth}
    \centering
    \IfFileExists{images/sql_data_types.png}{
      \includegraphics[width=0.98\linewidth,height=0.62\textheight,keepaspectratio]{images/sql_data_types.png}
    }{
      \fbox{\parbox[c][0.4\textheight][c]{0.8\linewidth}{\centering Sơ đồ nhóm kiểu dữ liệu SQL}}
    }
    \\[1mm]
    \tiny\textit{Hình 1: Cấu trúc phân nhóm kiểu dữ liệu phổ biến trong SQL}
  \end{column}
  \begin{column}{0.50\textwidth}
    \begin{block}{6 nhóm dữ liệu chính trong MySQL}
      \begin{itemize}\setlength{\itemsep}{2pt}
        \item \kw{Numeric}: Số nguyên (\texttt{INT}), số thực chính xác (\texttt{DECIMAL}), số gần đúng (\texttt{FLOAT}).
        \item \kw{String}: Chuỗi ngắn (\texttt{CHAR}, \texttt{VARCHAR}), văn bản dài (\texttt{TEXT}).
        \item \kw{Date/Time}: \texttt{DATE}, \texttt{TIME}, \texttt{DATETIME}, \texttt{TIMESTAMP}.
        \item \kw{Binary}: \texttt{BINARY}, \texttt{VARBINARY}, \texttt{BLOB}.
        \item \kw{Boolean}: \texttt{BOOLEAN} (bí danh cho \texttt{TINYINT(1)}).
        \item \kw{Special}: \texttt{JSON}, \texttt{GEOMETRY} (Spatial GIS).
      \end{itemize}
    \end{block}
  \end{column}
\end{columns}
\end{frame}

%------------------------------------------------
% Slide 6: Quy trinh 5 buoc lua chon
\begin{frame}{3. Quy trình 5 bước lựa chọn kiểu dữ liệu tối ưu}
\begin{center}
\begin{tikzpicture}[node distance=0.8cm, auto,
  stepnode/.style={rectangle, draw=dbblue!80, fill=dblight, text=black, rounded corners=3pt, font=\scriptsize, align=center, inner sep=3.5pt, text width=0.88\linewidth}]
  \node[stepnode] (s1) {\textbf{Bước 1: Xác định thực thể \& thông tin cần lưu} (Họ tên, mã đơn hàng, ngày sinh, đơn giá, số dư)};
  \node[stepnode, below=0.18cm of s1] (s2) {\textbf{Bước 2: Xác định bản chất ngữ nghĩa} (Số đếm được, văn bản tìm kiếm, thời điểm, cờ logic)};
  \node[stepnode, below=0.18cm of s2] (s3) {\textbf{Bước 3: Chọn nhóm kiểu phù hợp} (Numeric, String, Temporal, Binary, JSON)};
  \node[stepnode, below=0.18cm of s3] (s4) {\textbf{Bước 4: Xác định giới hạn \& độ chính xác} (\texttt{INT UNSIGNED}, \texttt{VARCHAR(100)}, \texttt{DECIMAL(10,2)})};
  \node[stepnode, below=0.18cm of s4] (s5) {\textbf{Bước 5: Bổ sung ràng buộc toàn vẹn} (\texttt{NOT NULL}, \texttt{CHECK}, \texttt{DEFAULT}, \texttt{UNIQUE})};
  \draw[-{Latex[length=2mm]}, thick, color=dbblue] (s1) -- (s2);
  \draw[-{Latex[length=2mm]}, thick, color=dbblue] (s2) -- (s3);
  \draw[-{Latex[length=2mm]}, thick, color=dbblue] (s3) -- (s4);
  \draw[-{Latex[length=2mm]}, thick, color=dbblue] (s4) -- (s5);
\end{tikzpicture}
\end{center}
\end{frame}

%------------------------------------------------
\section{Kiểu dữ liệu số (Numeric)}
% Slide 7: So nguyen
\begin{frame}{4. Kiểu số nguyên (Integer Types) trong MySQL}
\begin{center}
\resizebox{0.96\textwidth}{!}{
\begin{tabular}{lcccc}
\toprule
\textbf{Kiểu dữ liệu} & \textbf{Bộ nhớ} & \textbf{Miền giá trị (Signed)} & \textbf{Miền giá trị (Unsigned)} & \textbf{Ứng dụng tiêu biểu} \\
\midrule
\texttt{TINYINT} & 1 byte & $-128 \dots 127$ & $0 \dots 255$ & Trạng thái, cờ boolean, số thứ tự nhỏ \\
\texttt{SMALLINT} & 2 bytes & $-32{,}768 \dots 32{,}767$ & $0 \dots 65{,}535$ & Năm sản xuất, số phòng, số lượng nhỏ \\
\texttt{MEDIUMINT} & 3 bytes & $-8{,}388{,}608 \dots 8{,}388{,}607$ & $0 \dots 16{,}777{,}215$ & Mã danh mục, bảng tra cứu vừa \\
\texttt{INT} & 4 bytes & $\approx -2.14 \times 10^9 \dots 2.14 \times 10^9$ & $0 \dots \approx 4.29 \times 10^9$ & Khóa chính phổ biến, số lượng tồn kho \\
\texttt{BIGINT} & 8 bytes & $\approx -9.22 \times 10^{18} \dots 9.22 \times 10^{18}$ & $0 \dots \approx 1.84 \times 10^{19}$ & Log giao dịch hệ thống lớn, Big Data \\
\bottomrule
\end{tabular}
}
\end{center}

\smallgap
\begin{alertblock}{Nguyên tắc tối ưu hóa bộ nhớ}
  Tránh dùng \texttt{BIGINT} vô tội vạ cho mọi bảng! Một bảng 100 triệu dòng dùng \texttt{BIGINT} thay vì \texttt{INT} sẽ tiêu tốn thêm $\approx 400\text{ MB}$ chỉ cho riêng cột đó, làm chậm tốc độ quét index trên RAM.
\end{alertblock}
\end{frame}

%------------------------------------------------
% Slide 8: So thap phan DECIMAL vs FLOAT
\begin{frame}{5. Số thập phân chính xác vs Số thực gần đúng}
\begin{columns}[t]
  \begin{column}{0.48\textwidth}
    \begin{block}{\texttt{DECIMAL(P, D)} -- Số chính xác}
      \begin{itemize}\setlength{\itemsep}{1.5pt}
        \item Lưu trữ \kw{chính xác tuyệt đối} từng chữ số dưới dạng nhị phân đóng gói.
        \item $P$ (\textit{Precision}): Tổng chữ số (tối đa 65).
        \item $D$ (\textit{Scale}): Số chữ số thập phân (tối đa 30).
        \item \textbf{Ví dụ}: \texttt{DECIMAL(10,2)} $\rightarrow$ tối đa $99{,}999{,}999.99$.
        \item \kw{Bắt buộc cho}: Tiền tệ, lương, số dư tài khoản ngân hàng, hóa đơn.
      \end{itemize}
    \end{block}
  \end{column}
  \begin{column}{0.48\textwidth}
    \begin{block}{\texttt{FLOAT} \& \texttt{DOUBLE} -- Số gần đúng}
      \begin{itemize}\setlength{\itemsep}{1.5pt}
        \item Chuẩn IEEE 754 biểu diễn số thực.
        \item \texttt{FLOAT}: 4 bytes ($\approx 7$ chữ số tin cậy).
        \item \texttt{DOUBLE}: 8 bytes ($\approx 15$ chữ số tin cậy).
        \item Tính toán rất nhanh nhưng có \kw{sai số làm tròn (rounding error)}.
        \item \kw{Phù hợp cho}: Tọa độ GPS, cảm biến đo nhiệt độ, đo lường khoa học.
      \end{itemize}
    \end{block}
  \end{column}
\end{columns}

\begin{center}
  \footnotesize \textcolor{dborange}{\textbf{Cảnh báo vàng:}} Không bao giờ dùng \texttt{FLOAT} hay \texttt{DOUBLE} để lưu trữ tiền tệ!
\end{center}
\end{frame}

%------------------------------------------------
\section{Kiểu chuỗi (String)}
% Slide 9: CHAR vs VARCHAR
\begin{frame}{6. Kiểu chuỗi: \texttt{CHAR} vs \texttt{VARCHAR}}
\begin{columns}[t]
  \begin{column}{0.48\textwidth}
    \begin{block}{\texttt{CHAR(n)} -- Chuỗi cố định}
      \begin{itemize}\setlength{\itemsep}{1.5pt}
        \item Chiều dài luôn cố định là $n$ ký tự ($0 \le n \le 255$).
        \item Ngắn hơn $n$ sẽ được \kw{đệm dấu cách} vào cuối.
        \item Truy xuất trực tiếp nhanh hơn vì offset cố định.
        \item \kw{Ứng dụng}: Mã quốc gia (\texttt{'VN'}, \texttt{'US'}), mã tiền tệ (\texttt{'VND'}), mã bưu điện, hash MD5.
      \end{itemize}
    \end{block}
  \end{column}
  \begin{column}{0.48\textwidth}
    \begin{block}{\texttt{VARCHAR(n)} -- Chuỗi biến đổi}
      \begin{itemize}\setlength{\itemsep}{1.5pt}
        \item Chiều dài linh hoạt theo nội dung thực tế (tối đa $n$ ký tự).
        \item Tốn thêm \kw{1 hoặc 2 bytes tiền tố} lưu độ dài chuỗi.
        \item Tiết kiệm dung lượng đĩa khi dữ liệu có độ dài chênh lệch lớn.
        \item \kw{Ứng dụng}: Họ tên, email, địa chỉ, tiêu đề.
      \end{itemize}
    \end{block}
  \end{column}
\end{columns}

\smallgap
\begin{exampleblock}{So sánh lưu trữ giá trị \texttt{'AB'} (với bảng mã 1-byte)}
  \begin{itemize}\setlength{\itemsep}{1pt}
    \item \texttt{CHAR(5)}: Chiếm 5 bytes (2 ký tự \texttt{'AB'} + 3 ký tự khoảng trắng đệm).
    \item \texttt{VARCHAR(5)}: Chiếm 3 bytes (2 ký tự \texttt{'AB'} + 1 byte lưu độ dài $L=2$).
  \end{itemize}
\end{exampleblock}
\end{frame}

%------------------------------------------------
% Slide 10: TEXT va utf8mb4
\begin{frame}{7. Chuỗi văn bản dài (\texttt{TEXT}) \& Bộ mã \texttt{utf8mb4}}
\begin{columns}[t]
  \begin{column}{0.48\textwidth}
    \begin{block}{Các kiểu văn bản lớn (\texttt{TEXT})}
      \begin{itemize}\setlength{\itemsep}{1.5pt}
        \item \texttt{TINYTEXT}: Tối đa 255 bytes.
        \item \texttt{TEXT}: Tối đa $64\text{ KB}$ (mô tả sản phẩm).
        \item \texttt{MEDIUMTEXT}: Tối đa $16\text{ MB}$ (nội dung bài báo).
        \item \texttt{LONGTEXT}: Tối đa $4\text{ GB}$ (tài liệu đồ sộ).
        \item Lưu trữ ngoài dòng (off-page) trong InnoDB khi nội dung lớn, tránh phân mảnh bảng.
      \end{itemize}
    \end{block}
  \end{column}
  \begin{column}{0.48\textwidth}
    \begin{block}{Charset \texttt{utf8mb4} \& Collation}
      \begin{itemize}\setlength{\itemsep}{1.5pt}
        \item \kw{utf8mb4} là chuẩn UTF-8 đầy đủ của MySQL (1--4 bytes/ký tự).
        \item Hỗ trợ toàn diện tiếng Việt có dấu, chữ tượng hình và biểu tượng cảm xúc.
        \item \texttt{utf8mb4\_0900\_ai\_ci}: Collation chuẩn từ MySQL 8.0, sắp xếp theo Unicode 9.0, không phân biệt hoa thường và không phân biệt dấu khi so sánh.
      \end{itemize}
    \end{block}
  \end{column}
\end{columns}
\end{frame}

%------------------------------------------------
\section{Kiểu ngày giờ (Date \& Time)}
% Slide 11: DATETIME vs TIMESTAMP
\begin{frame}{8. Kiểu ngày giờ: \texttt{DATETIME} vs \texttt{TIMESTAMP}}
\begin{center}
\resizebox{0.96\textwidth}{!}{
\begin{tabular}{lcp{4.5cm}p{5.5cm}}
\toprule
\textbf{Kiểu} & \textbf{Bộ nhớ} & \textbf{Miền giá trị hợp lệ} & \textbf{Đặc tính Timezone \& Ứng dụng} \\
\midrule
\texttt{DATETIME} & 5 bytes & \texttt{1000-01-01} $\dots$ \texttt{9999-12-31} & \textbf{Cố định}: Không phụ thuộc múi giờ session. Phù hợp cho ngày sinh, hạn bảo hành, lịch hẹn. \\
\texttt{TIMESTAMP} & 4 bytes & \texttt{1970-01-01} $\dots$ \texttt{2038-01-19} UTC & \textbf{Tự động đổi về UTC}: Đổi theo múi giờ client khi đọc. Thích hợp cho audit log, \texttt{created\_at}. \\
\texttt{DATE} & 3 bytes & \texttt{1000-01-01} $\dots$ \texttt{9999-12-31} & Chỉ lưu năm-tháng-ngày (\texttt{YYYY-MM-DD}). Tiết kiệm cho ngày sinh, ngày hóa đơn. \\
\texttt{TIME} & 3 bytes & \texttt{-838:59:59} $\dots$ \texttt{838:59:59} & Lưu giờ trong ngày hoặc khoảng thời gian trôi qua. \\
\bottomrule
\end{tabular}
}
\end{center}

\smallgap
\begin{alertblock}{Cảnh báo giới hạn năm 2038 của TIMESTAMP}
  Do \texttt{TIMESTAMP} chỉ dùng số nguyên có dấu 32-bit tính từ epoch 1970, nó sẽ tràn vào ngày 19/01/2038. Đối với các mốc thời gian trong tương lai xa, hãy sử dụng \kw{\texttt{DATETIME}}!
\end{alertblock}
\end{frame}

%------------------------------------------------
\section{Kiểu đặc biệt \& Thực hành}
% Slide 12: Binary, Boolean, JSON
\begin{frame}{9. Kiểu Binary, Boolean \& JSON hiện đại}
\begin{columns}[t]
  \begin{column}{0.48\textwidth}
    \begin{block}{\texttt{BLOB} vs \texttt{VARBINARY}}
      \begin{itemize}\setlength{\itemsep}{1.5pt}
        \item \texttt{BINARY} / \texttt{VARBINARY}: Chuỗi byte thô cố định hoặc linh hoạt (lưu hash mật khẩu, UUID binary 16 bytes).
        \item \texttt{BLOB}: Lưu file nhị phân (ảnh, file đính kèm).
        \item \textit{Thực tế tốt nhất}: Nên lưu file media trên Cloud Storage (S3, CDN) và chỉ lưu URL chuỗi trong database.
      \end{itemize}
    \end{block}
  \end{column}
  \begin{column}{0.48\textwidth}
    \begin{block}{\texttt{BOOLEAN} \& \texttt{JSON}}
      \begin{itemize}\setlength{\itemsep}{1.5pt}
        \item Trong MySQL, \texttt{BOOLEAN} là alias của \texttt{TINYINT(1)} ($0 = \text{FALSE}$, $1 = \text{TRUE}$).
        \item \texttt{JSON}: Kiểu dữ liệu gốc (MySQL 5.7+). Tự động validate cú pháp JSON, hỗ trợ index trên virtual column và truy vấn nhanh qua toán tử \texttt{->} và \texttt{->>}.
      \end{itemize}
    \end{block}
  \end{column}
\end{columns}
\end{frame}

%------------------------------------------------
% Slide 13: Demo code CREATE TABLE
\begin{frame}[fragile]{10. Thực hành: Khai báo đa dạng kiểu dữ liệu trong bảng}
\begin{lstlisting}
CREATE TABLE products (
    product_id   INT UNSIGNED AUTO_INCREMENT PRIMARY KEY, -- Khoa chinh so nguyen duong
    product_code CHAR(8) NOT NULL UNIQUE,                -- Ma san pham co dinh 8 ky tu
    product_name VARCHAR(120) NOT NULL,                  -- Ten san pham linh hoat
    category_id  SMALLINT UNSIGNED NOT NULL,             -- Ma danh muc < 65,535
    cost_price   DECIMAL(10,2) NOT NULL,                 -- Gia von: chinh xac tuyet doi
    sale_price   DECIMAL(10,2) NOT NULL,                 -- Gia ban: chinh xac tuyet doi
    weight_kg    FLOAT(5,2),                             -- Can nang: chap nhan sai so nho
    description  TEXT,                                   -- Mo ta bai viet dai
    attributes   JSON,                                   -- Thong so ky thuat dong
    is_active    BOOLEAN NOT NULL DEFAULT TRUE,          -- Con kinh doanh (TINYINT(1))
    created_at   DATETIME NOT NULL DEFAULT CURRENT_TIMESTAMP -- Moc thoi gian tao
) ENGINE=InnoDB DEFAULT CHARSET=utf8mb4 COLLATE=utf8mb4_0900_ai_ci;
\end{lstlisting}
\end{frame}

%------------------------------------------------
% Slide 14: Sai lam thuong gap
\begin{frame}{11. Các sai lầm kinh điển khi chọn kiểu dữ liệu}
\begin{columns}[t]
  \begin{column}{0.48\textwidth}
    \begin{alertblock}{Sai lầm thiết kế phổ biến}
      \begin{itemize}\setlength{\itemsep}{2pt}
        \item \textbf{Lưu số điện thoại bằng INT}: Bị mất số \texttt{0} đầu (\texttt{0912...} thành \texttt{912...}), mất dấu \texttt{+}. Luôn dùng \texttt{VARCHAR(20)}.
        \item \textbf{Lưu tiền bằng FLOAT}: Dễ gặp sai số làm tròn số học ($0.1 + 0.2 \ne 0.3$). Luôn dùng \texttt{DECIMAL}.
        \item \textbf{Lưu ngày tháng bằng VARCHAR}: Không so sánh ngày chuẩn, không tính toán khoảng cách ngày. Dùng \texttt{DATE}.
      \end{itemize}
    \end{alertblock}
  \end{column}
  \begin{column}{0.48\textwidth}
    \begin{alertblock}{Lạm dụng và lãng phí}
      \begin{itemize}\setlength{\itemsep}{2pt}
        \item \textbf{Lạm dụng VARCHAR(255) mọi nơi}: Tốn bộ nhớ RAM nội bộ khi sắp xếp \texttt{ORDER BY} hoặc tạo bảng tạm.
        \item \textbf{Dùng BIGINT cho cờ nhị phân}: Lãng phí bộ nhớ gấp 8 lần so với \texttt{TINYINT} hoặc \texttt{BOOLEAN}.
        \item \textbf{Lưu video/file to vào BLOB}: Gây phình to kích thước database, chậm backup và restore hệ thống.
      \end{itemize}
    \end{alertblock}
  \end{column}
\end{columns}
\end{frame}

%------------------------------------------------
% Slide 15: Quiz nhanh cung co
\begin{frame}{12. Câu hỏi kiểm tra nhanh (Quick Quiz)}
\begin{columns}[t]
  \begin{column}{0.48\textwidth}
    \begin{block}{Câu 1: Chọn kiểu cho đơn giá}
      Cột lưu đơn giá bán hàng nên khai báo kiểu nào?
      \begin{itemize}
        \item A. \texttt{FLOAT}
        \item B. \textbf{\kw{DECIMAL(10,2)}} (Chính xác!)
        \item C. \texttt{VARCHAR(50)}
        \item D. \texttt{BIGINT}
      \end{itemize}
    \end{block}
  \end{column}
  \begin{column}{0.48\textwidth}
    \begin{block}{Câu 2: Số điện thoại khách hàng}
      Vì sao nên dùng \texttt{VARCHAR} cho số điện thoại?
      \begin{itemize}
        \item A. Số điện thoại không dùng tính toán.
        \item B. Giữ nguyên số \texttt{0} đầu và dấu \texttt{+}.
        \item C. Tránh nguy cơ tràn số nguyên.
        \item D. \textbf{\kw{Cả A, B và C đều đúng}}.
      \end{itemize}
    \end{block}
  \end{column}
\end{columns}
\smallgap
\begin{exampleblock}{Câu 3: Phân biệt CHAR và VARCHAR}
  \textbf{Hỏi}: Khi lưu mã quốc gia 2 ký tự (\texttt{'VN'}, \texttt{'JP'}, \texttt{'US'}), nên chọn \texttt{CHAR(2)} hay \texttt{VARCHAR(2)}?\\
  \textbf{Trả lời}: Chọn \kw{\texttt{CHAR(2)}} vì độ dài cố định, truy xuất trực tiếp, không tốn thêm byte lưu tiền tố chiều dài!
\end{exampleblock}
\end{frame}

%------------------------------------------------
% Slide 16: Bai tap ve nha
\begin{frame}{13. Bài tập vận dụng thiết kế bảng}
\begin{block}{Bài tập thực hành tại lớp}
  Hãy viết câu lệnh \texttt{CREATE TABLE} cho bảng \texttt{student\_profiles} để quản lý hồ sơ sinh viên với các yêu cầu:
  \begin{enumerate}\setlength{\itemsep}{2pt}
    \item \textbf{Mã sinh viên}: Số nguyên dương tự tăng, làm khóa chính.
    \item \textbf{Mã định danh SV}: Chuỗi đúng 10 ký tự (\texttt{CHAR(10)}), không trùng lặp.
    \item \textbf{Họ và tên}: Tối đa 100 ký tự có dấu tiếng Việt (\texttt{VARCHAR(100)}).
    \item \textbf{Ngày sinh}: Chỉ cần lưu ngày tháng năm (\texttt{DATE}).
    \item \textbf{Điểm trung bình (GPA)}: Thang điểm 4 hoặc 10, có 2 chữ số thập phân (\texttt{DECIMAL(4,2)}).
    \item \textbf{Tình trạng nợ học phí}: Giá trị logic Đúng / Sai (\texttt{BOOLEAN}).
    \item \textbf{Thời điểm nộp hồ sơ}: Tự động ghi nhận mốc thời gian hệ thống (\texttt{DATETIME}).
  \end{enumerate}
\end{block}
\end{frame}

%------------------------------------------------
% Slide 17: Tong ket
\begin{frame}{Tổng kết bài học}
\begin{columns}[t]
  \begin{column}{0.5\textwidth}
    \begin{block}{Ba nguyên tắc vàng ghi nhớ}
      \begin{itemize}\setlength{\itemsep}{2pt}
        \item \kw{Đúng bản chất}: Số dùng để tính, chuỗi dùng để đọc, thời gian dùng để đo lường.
        \item \kw{Vừa đủ phạm vi}: Tránh chọn quá nhỏ gây tràn dữ liệu, tránh chọn quá to gây lãng phí RAM.
        \item \kw{Chính xác tuyệt đối}: Dữ liệu tài chính, tiền tệ bắt buộc dùng \texttt{DECIMAL}.
      \end{itemize}
    \end{block}
  \end{column}
  \begin{column}{0.46\textwidth}
    \begin{exampleblock}{Tài liệu học tập tiếp theo}
      \begin{itemize}\setlength{\itemsep}{2pt}
        \item Bài tiếp theo: \textbf{CREATE TABLE \& Ràng buộc toàn vẹn} (Primary Key, Foreign Key, Check, Unique).
        \item Lab thực hành: Kiểm tra dung lượng trên đĩa của các kiểu dữ liệu qua \texttt{information\_schema}.
      \end{itemize}
    \end{exampleblock}
  \end{column}
\end{columns}
\end{frame}

\end{document}
